\section*{अध्याय \ref{ch:b2:conjugate-additions} --- संयुग्मित कार्बोनिल: माइकल और रॉबिन्सन}

\begin{solution}{exo:b2:conjugate-additions:1}
\begin{itemize}
  \item \ce{CH3MgBr}: मुख्यतः 1,2 (कठोर, अनुत्क्रमणीय), 1-मेथिलसाइक्लोहेक्स-2-ईन-1-ऑल।
  \item \ce{(CH3)2CuLi}: 1,4, 3-मेथिलसाइक्लोहेक्सेनोन।
  \item क्षारक के साथ थायोफ़ीनॉल: 1,4 (थायोलेट मृदु है), 3-(फ़ेनिलसल्फ़ेनिल)साइक्लोहेक्सेनोन।
  \item \ce{LiAlH4}: 1,2, साइक्लोहेक्स-2-ईन-1-ऑल।
\end{itemize}
\end{solution}

\begin{solution}{exo:b2:conjugate-additions:2}
दाता: पेंटेन-2,4-डाइओन, नाइट्रोमेथेन, डाइएथिल प्रोपेनडाइओएट (अम्लीय \ce{C-H})। ग्राही:
प्रोपीननाइट्राइल, मेथिल प्रोपीनोएट, ब्यूट-3-ईन-2-ओन (किसी ग्राही समूह से संयुग्मित \ce{C=C})।
\end{solution}

\begin{solution}{exo:b2:conjugate-additions:3}
\ce{C4H9Br + 2Li -> C4H9Li + LiBr}, फिर \ce{2C4H9Li + CuI -> (C4H9)2CuLi + LiI}, ईथर में निम्न
ताप पर, अक्रिय गैस के अंतर्गत।
\end{solution}

\begin{solution}{exo:b2:conjugate-additions:4}\emergencystretch=8em
डाइएथिल (3-ऑक्सोब्यूटिल)प्रोपेनडाइओएट, \ce{(EtOOC)2CHCH2CH2COCH3}। जल-अपघटन प्रतिस्थापित
प्रोपेनडाइओइक अम्ल देता है, जो गर्म करने पर \ce{CO2} खो देता है (\cref{prop:b2:enolates-aldol:decarboxylation}):
5-ऑक्सोहेक्सेनोइक अम्ल, \ce{HOOCCH2CH2CH2COCH3}।
\end{solution}

\begin{solution}{exo:b2:conjugate-additions:5}
निम्न ताप पर ईथर में लिथियम डाइमेथिलक्यूप्रेट, फिर जलीय उपचार-पश्चात् प्रक्रिया: \omterm{def:b2:conjugate-additions:addition-modes}{संयुग्मी योग}।
मेथिलमैग्नीशियम ब्रोमाइड, अधिक कठोर, अधिकांशतः 1,2 जुड़ता है और मुख्य उत्पाद के रूप में
1-मेथिलसाइक्लोहेक्स-2-ईन-1-ऑल देता है।
\end{solution}

\begin{solution}{exo:b2:conjugate-additions:6}
\ce{Et3N} थोड़े \ce{EtSH} से प्रोटॉन हटाकर \ce{EtS-} देता है; थायोलेट $\beta$ कार्बन से जुड़ता है,
बना \omterm{def:b2:enolates-aldol:enolate}{ईनोलेट} किसी अन्य \ce{EtSH} का प्रोटॉन ले लेता है, जिससे \ce{EtS-} का पुनर्जनन होता है। उत्पाद:
4-(एथिलसल्फ़ेनिल)ब्यूटेन-2-ओन। सल्फ़र बड़ा और ध्रुवणीय है, उसका एकाकी युग्म ऊँची ऊर्जा पर है: एक \omterm{def:b2:frontier-orbitals:hard-soft}{मृदु
नाभिकरागी}, \omterm{def:b2:frontier-orbitals:control}{कक्षक नियंत्रण} में, अतः वह वहाँ आक्रमण करता है जहाँ \omterm{def:b2:frontier-orbitals:frontier}{LUMO} सबसे बड़ा है, और योग
उत्क्रमणीय है, जो अधिक स्थायी 1,4-योगज के अनुकूल भी है।
\end{solution}

\begin{solution}{exo:b2:conjugate-additions:7}
सायनाइड अधिक धनात्मक कार्बोनिल कार्बन पर तेज़ी से आक्रमण करता है: सायनोहाइड्रिन गतिक उत्पाद है। उसका
बनना उत्क्रमणीय है; उच्चतर ताप पर, समय मिलने पर, सायनाइड मुक्त होकर
$\beta$ कार्बन पर फिर जुड़ता है, और वह योगज देता है जो प्रबलतर \ce{C=O} $\pi$ आबंध बनाए रखता है, अर्थात् ऊष्मागतिक उत्पाद।
\end{solution}

\begin{solution}{exo:b2:conjugate-additions:8}
2-मेथिलसाइक्लोहेक्सेन-1,3-डाइओन के साथ: वीलैंड--मीशर कीटोन (\omterm{def:b2:conjugate-additions:robinson}{रॉबिन्सन वलयन} का चित्र)।
साइक्लोहेक्सेनोन के साथ: 4,4a,5,6,7,8-हेक्साहाइड्रोनैफ़्थलीन-2(3\textit{H})-ओन, एक द्विचक्रीय ईनोन जिसका
\ce{C=C} एक वलय-संगलन कार्बन पर समाप्त होता है, बिना किसी कोणीय मेथिल समूह के।
\end{solution}

\begin{solution}{exo:b2:conjugate-additions:9}
\ce{CH3COCH2CH2CH2COCH3}: C2 और C6 कार्बोनिल हैं, C2 से C6 तक दोनों को सम्मिलित करते हुए पाँच कार्बन,
अर्थात् एक 1,5-डाइकीटोन। C3--C4 काटिए: दाता, C3 पर प्रोपेनोन का \omterm{def:b2:enolates-aldol:enolate}{ईनोलेट} (बेहतर, एथिल 3-ऑक्सोब्यूटेनोएट, जिसका
एस्टर बाद में जल-अपघटन और \omterm{def:b2:enolates-aldol:decarboxylation}{विकार्बोक्सिलन} से हटा दिया जाता है); ग्राही, ब्यूट-3-ईन-2-ओन। परिस्थितियाँ:
एथिल 3-ऑक्सोब्यूटेनोएट, ब्यूट-3-ईन-2-ओन, एथेनॉल में उत्प्रेरकी सोडियम एथॉक्साइड; फिर जलीय अम्ल और
ताप।
\end{solution}

\begin{solution}{exo:b2:conjugate-additions:10}
\omterm{def:b2:enolates-aldol:aldol}{ऐल्डोल} एक \ce{C-C} आबंध बनाता है परंतु एक \ce{C=O} $\pi$ आबंध को \ce{C-O} $\sigma$ आबंध में बदलता है;
संतुलन छोटा है और दो अणु एक बनते हैं, अतः $K$ साधारण है और प्रतीप-ऐल्डोल आसान है।
\omterm{def:b2:conjugate-additions:michael}{माइकल योग} एक \ce{C-C} आबंध बनाता है, एक दुर्बलतर \ce{C=C} $\pi$ आबंध की क़ीमत पर, और दोनों
\ce{C=O} बनाए रखता है: यह स्पष्ट रूप से अनुकूल है। उलटी अभिक्रिया के लिए कीटोन \omterm{def:b2:enolates-aldol:enolate}{ईनोलेट} को स्थायीकृत प्रोपेनडाइओएट ऋणायन बाहर निकालना पड़ता;
साम्य योगज की ओर इतना झुका है कि उत्प्रेरकी क्षारक के अंतर्गत योगज बना रहता है।
\end{solution}

\begin{solution}{exo:b2:conjugate-additions:11}
\ce{C(COOEt)2(CH2CH2COOCH3)2}। पहले \omterm{def:b2:conjugate-additions:michael}{माइकल योग} के बाद केंद्रीय कार्बन पर दो एस्टरों के बीच अभी भी एक
हाइड्रोजन है, जो पहले जितना ही अम्लीय है: क्षारक उसे हटाता है और दूसरा योग होता है।
\end{solution}

\begin{solution}{exo:b2:conjugate-additions:12}
अधिक प्रतिस्थापित \omterm{def:b2:enolates-aldol:enolate}{ईनोलेट} (मेथिल धारण करने वाले कार्बन की ओर) से होकर: वह ऑक्टैलोन जिसमें मेथिल
समूह वलय-संगलन कार्बन पर है, 4a-मेथिल-4,4a,5,6,7,8-हेक्साहाइड्रोनैफ़्थलीन-2(3\textit{H})-ओन। कम
प्रतिस्थापित \omterm{def:b2:enolates-aldol:enolate}{ईनोलेट} (\ce{CH2} वाली ओर) से होकर: मेथिल समूह संगलन के बगल वाले किसी वलय कार्बन पर पहुँचता है।
पहला साम्यकारी परिस्थितियों में अनुकूल है (अपने ऐल्कोहॉल में कोई ऐल्कॉक्साइड, या किसी \omterm{def:b2:amines:class}{द्वितीयक ऐमीन} और अम्ल के साथ
\omterm{def:b2:amines:imine}{एनामीन} मार्ग), जो ऊष्मागतिक, अधिक प्रतिस्थापित \omterm{def:b2:enolates-aldol:enolate}{ईनोलेट} देती हैं
(\cref{prop:b2:enolates-aldol:enolate-control})।
\end{solution}

\begin{solution}{pb:b2:conjugate-additions:1}
\textbf{1.} एक साइक्लोहेक्सेन वलय जिसके C1 और C3 पर \ce{C=O} और C2 पर एक मेथिल समूह है; C2 पर स्थित हाइड्रोजन
सबसे अम्लीय है।
\textbf{2.} इसे हटाने से ऐसा ऋणायन मिलता है जिसका आवेश C2 और दोनों ऑक्सीजनों पर विस्थानीकृत है;
साइक्लोहेक्सेनोन में केवल एक ऑक्सीजन उसे साझा करता है, और अम्लता इतनी दुर्बल है कि जल में मापी नहीं जा सकती।
\textbf{3.} C2 पर आवेश; \ce{C1=C2}, साथ में C1 पर \ce{O^-}; \ce{C2=C3}, साथ में C3 पर \ce{O^-}।
\textbf{4.} हाँ: हाइड्रॉक्साइड के साथ $K = 10^{14.0 - 5.3} = 10^{8.7}$ (संयुग्मी अम्ल के रूप में जल,
$K_e$ से); एथॉक्साइड के साथ $10^{15.9-5.3} = 10^{10.6}$। दोनों में से कोई भी पूर्ण है।
% ledger: ke:25C, pka:ethanol
\textbf{5.} \ce{C1(OH)=C2(CH3)}: C2 पर अभी भी एक हाइड्रोजन है, जो \omterm{def:b2:enolates-aldol:tautomerism}{ईनॉल} के लिए पर्याप्त है; मेथिल
समूह केवल दूसरे का स्थान लेता है।
\textbf{6.} मृदु: आवेश तीन परमाणुओं पर फैला है और ऋणायन स्थायीकृत है। यह
$\beta$ कार्बन पर आक्रमण करता है (\cref{prop:b2:conjugate-additions:regio})।
\textbf{7.} \omterm{def:b2:enolates-aldol:enolate}{ईनोलेट} का C2 ब्यूट-3-ईन-2-ओन के \ce{CH2} सिरे से आबंधित होता है; बना \omterm{def:b2:enolates-aldol:enolate}{ईनोलेट} एक
प्रोटॉन लेता है: 2-मेथिल-2-(3-ऑक्सोब्यूटिल)साइक्लोहेक्सेन-1,3-डाइओन।
\textbf{8.} हर नया \omterm{def:b2:enolates-aldol:enolate}{ईनोलेट} डाइओन के किसी अणु से (या विलायक से) एक प्रोटॉन लेता है और इस प्रकार
दाता \omterm{def:b2:enolates-aldol:enolate}{ईनोलेट} का पुनर्जनन करता है: क्षारक उपभुक्त नहीं होता।
\textbf{9.} \omterm{def:b2:enolates-aldol:enolate}{ईनोलेट} उभयदंती है; किसी मृदु कार्बन इलेक्ट्रॉनरागी के साथ, \omterm{def:b2:frontier-orbitals:control}{कक्षक नियंत्रण} में, यह कार्बन पर
अभिक्रिया करता है, और \ce{C-C} योगज दो \ce{C=O} आबंध बनाए रखता है, जो प्रबलतर संयोजन है; कोई
\ce{O}-योगज, यदि बने, तो वापस लौट जाता।
\textbf{10.} एक: C2, जो C1, C3, मेथिल और शृंखला से आबंधित है।
\textbf{11.} शृंखला का कार्बोनिल कार्बन और वलय का कार्बोनिल C1 (या C3): \ce{C(=O)}--\ce{CH2}--\ce{CH2}--C2--C1,
दोनों सिरों को सम्मिलित करते हुए पाँच कार्बन।
\textbf{12.} पूरे डाइओन का, जो अधिक मूल्यवान अभिकर्मक है, उपभोग करने के लिए; केवल हल्का, क्योंकि ब्यूट-3-ईन-2-ओन
अत्यंत विषैला है और बहुलकित हो जाता है।
\textbf{13.} वलय के C4 और C6 (C3 और C1 के बगल वाले), शृंखला का अपने कार्बोनिल के बगल वाला \ce{CH2}, और शृंखला का
अंतस्थ \ce{CH3}। C2 पर कोई नहीं।
\textbf{14.} अंतस्थ \ce{CH3} का \omterm{def:b2:enolates-aldol:enolate}{ईनोलेट} वलय के एक कार्बोनिल पर आक्रमण करता है: बंद हुए वलय में वह
कार्बन, शृंखला का कार्बोनिल कार्बन, दो \ce{CH2}, C2 और आक्रांत कार्बोनिल कार्बन गिने जाते हैं: छह परमाणु।
\textbf{15.} शृंखला के भीतरी \ce{CH2} का \omterm{def:b2:enolates-aldol:enolate}{ईनोलेट} एक चतुःसदस्यी, तनावग्रस्त वलय बंद करता। शृंखला कार्बोनिल पर आक्रमण करता
वलय का कोई \omterm{def:b2:enolates-aldol:enolate}{ईनोलेट} (C4 या C6) भी षट्सदस्यी वलय बंद करता है, परंतु सेतुबद्ध: वह \omterm{def:b2:enolates-aldol:aldol}{ऐल्डोल}
निर्जलित नहीं हो सकता (द्वि-आबंध किसी छोटे द्विचक्रीय निकाय के सेतुशीर्ष पर पड़ता) और वापस लौट जाता है,
जबकि संगलित \omterm{def:b2:enolates-aldol:aldol}{ऐल्डोल} निर्जलीकरण द्वारा खपता रहता है।
\textbf{16.} \omterm{def:b2:enolates-aldol:aldol}{ऐल्डोल} का \ce{OH} भूतपूर्व वलय कार्बोनिल कार्बन पर है, जो अब वलय-संगलन पर है; भूतपूर्व
\ce{CH3} कार्बन (अब शृंखला कीटोन के सापेक्ष $\alpha$) के एक हाइड्रोजन के साथ जल की हानि
उस कीटोन से संयुग्मित \ce{C=C} देती है।
\textbf{17.} नया द्वि-आबंध कार्बोनिल से संयुग्मित है: \omterm{def:b2:enolates-aldol:enolate}{ईनोलेट} से होकर एक E1cB विलोपन,
जिसे तापन बढ़ावा देता है (\cref{prop:b2:enolates-aldol:aldol-equilibrium})।
\textbf{18.} एक संतृप्त कीटोन और एक $\alpha,\beta$-असंतृप्त कीटोन (एक साइक्लोहेक्सीनोन)।
\textbf{19.} मेथिल समूह धारण करने वाला वलय-संगलन कार्बन: चार भिन्न प्रतिस्थापी।
\textbf{20.} नहीं: डाइओन अकिरैल है, और उसके दोनों कार्बोनिलों पर समान रूप से आक्रमण होता है (शृंखला के सापेक्ष वे
एक-दूसरे के दर्पण प्रतिबिंब हैं); अकिरैल अभिकर्मकों के साथ उत्पाद रेसेमिक है। किरैल ऐमीन
उत्प्रेरक एक कार्बोनिल के अनुकूल होते हैं और एक प्रतिबिंबरूप आधिक्य में देते हैं (तृतीय वर्ष का खंड)।
\textbf{21.} डाइओन \ce{C7H10O2}: \qty{126.155}{g/mol}; \ce{C4H6O}: \qty{70.091}{g/mol}; \ce{C11H14O2}:
\qty{178.231}{g/mol}।
\textbf{22.} \ce{C7H10O2 + C4H6O -> C11H14O2 + H2O}: हर ओर C 11, H 16, O 3।
\textbf{23.} $10.0/126.155 = \qty{0.07927}{mol} = \qty{79.3}{mmol}$।
\textbf{24.} $0.07927 \times 178.231 \times 0.70 = \boldsymbol{\approx \qty{9.89}{g}}$
वीलैंड--मीशर कीटोन।
\end{solution}
